% GENERATED FILE. Edit canonical lessons, then run npm run generate:books.
% canonical-source-hash: fnv1a64:2060ddd77507cff9
% canonical-lessons: JA-C438-ganko, JA-C438-sekkachi, JA-C438-tanki, JA-C438-nonbiri, JA-C438-ibaru

\chapter{Stubborn and Hasty}
\label{ch:ja-stubborn-and-hasty}
\hlchaptermodality{438}

\begin{chapteropening}
\textbf{By the end of this chapter:} \emph{I can describe a stubborn or an impatient person.}

Complete the last lesson of chapter 438: I can describe a stubborn or an impatient person.
\end{chapteropening}

\section[\texorpdfstring{\ja{がんこ} (ganko)}{がんこ (ganko)}]{\ja{がんこ} (ganko) — stubborn; obstinate}

\label{lesson:JA-C438-ganko}



\begin{quote}
Before the new one: say the Japanese for modest, then the Japanese for spiteful.
\end{quote}

\subsection*{You'll want to know: \ja{がんこ}}
\textbf{\ja{がんこ}} — \emph{ganko} — \textquotedblleft{}stubborn; obstinate\textquotedblright{}.

\begin{quote}
\textbf{\ja{そふは} \ja{すこし} \ja{がんこです。}} — \emph{sofu wa sukoshi ganko desu} — \textquotedblleft{}My grandfather is a little stubborn.\textquotedblright{}
\end{quote}

A na-adjective: \ja{な} before a noun.

\begin{cousinweb}
\textbf{\ja{がんこ}} (\emph{ganko}) is Sino-Japanese (kango): two Chinese characters read with their on'yomi, readings borrowed from Chinese. They are \emph{gan}, \textquotedblleft{}stubborn\textquotedblright{} (Middle Chinese \emph{ngwaen}) and \emph{ko}, \textquotedblleft{}harden\textquotedblright{} (\emph{kuH}) (KANJIDIC2; Middle Chinese in Baxter's transcription).
\end{cousinweb}

\begin{grammarlens}[title={What you've built}]
One more word for describing what you lived and felt.
\end{grammarlens}

\subsection*{Your turn}
\begin{itemize}
\raggedright
  \item \emph{Say it:} \emph{ganko}
  \item \emph{Say it:} \emph{ganko}, once more
  \item \emph{Say it:} \emph{sofu wa sukoshi ganko desu} — My grandfather is a little stubborn.
  \item \emph{Recall:} say the Japanese for modest, then the Japanese for spiteful, then say \emph{ganko} again
\end{itemize}

\begin{tcolorbox}[breakable,colback=teal!4,colframe=teal!35!black,title={Before you move on}]
What does \ja{がんこ} mean? (\textquotedblleft{}Stubborn; obstinate\textquotedblright{}.) Say it once more.
\end{tcolorbox}
\section[\texorpdfstring{\ja{せっかち} (sekkachi)}{せっかち (sekkachi)}]{\ja{せっかち} (sekkachi) — impatient; hasty}

\label{lesson:JA-C438-sekkachi}



\begin{quote}
Before the new one: say the Japanese for spiteful, then the Japanese for stubborn.
\end{quote}

\subsection*{You'll want to know: \ja{せっかち}}
\textbf{\ja{せっかち}} — \emph{sekkachi} — \textquotedblleft{}impatient; hasty\textquotedblright{}.

\begin{quote}
\textbf{\ja{ちちは} \ja{せっかちです。}} — \emph{chichi wa sekkachi desu} — \textquotedblleft{}My father is impatient.\textquotedblright{}
\end{quote}

A na-adjective: \ja{な} before a noun.

\begin{cousinweb}
\textbf{\ja{せっかち}} (\emph{sekkachi}) is a native Japanese word (wago), inherited rather than borrowed from Chinese; JMdict writes it in kana (JMdict).
\end{cousinweb}

\begin{grammarlens}[title={What you've built}]
One more word for describing what you lived and felt.
\end{grammarlens}

\subsection*{Your turn}
\begin{itemize}
\raggedright
  \item \emph{Say it:} \emph{sekkachi}
  \item \emph{Say it:} \emph{sekkachi}, once more
  \item \emph{Say it:} \emph{chichi wa sekkachi desu} — My father is impatient.
  \item \emph{Recall:} say the Japanese for spiteful, then the Japanese for stubborn, then say \emph{sekkachi} again
\end{itemize}

\begin{tcolorbox}[breakable,colback=teal!4,colframe=teal!35!black,title={Before you move on}]
What does \ja{せっかち} mean? (\textquotedblleft{}Impatient; hasty\textquotedblright{}.) Say it once more.
\end{tcolorbox}
\section[\texorpdfstring{\ja{たんき} (tanki)}{たんき (tanki)}]{\ja{たんき} (tanki) — short-tempered}

\label{lesson:JA-C438-tanki}



\begin{quote}
Before the new one: say the Japanese for stubborn, then the Japanese for impatient.
\end{quote}

\subsection*{You'll want to know: \ja{たんき}}
\textbf{\ja{たんき}} — \emph{tanki} — \textquotedblleft{}short-tempered\textquotedblright{}.

\begin{quote}
\textbf{\ja{あには} \ja{すこし} \ja{たんきです。}} — \emph{ani wa sukoshi tanki desu} — \textquotedblleft{}My elder brother is a little short-tempered.\textquotedblright{}
\end{quote}

A na-adjective: \ja{な} before a noun.

\begin{cousinweb}
\textbf{\ja{たんき}} (\emph{tanki}) is Sino-Japanese (kango): two Chinese characters read with their on'yomi, readings borrowed from Chinese. They are \emph{tan}, \textquotedblleft{}short\textquotedblright{} (Middle Chinese \emph{twanX}) and \emph{ki}, \textquotedblleft{}spirit\textquotedblright{} (\emph{khj+jH}) (KANJIDIC2; Middle Chinese in Baxter's transcription).
\end{cousinweb}

\begin{grammarlens}[title={What you've built}]
One more word for describing what you lived and felt.
\end{grammarlens}

\subsection*{Your turn}
\begin{itemize}
\raggedright
  \item \emph{Say it:} \emph{tanki}
  \item \emph{Say it:} \emph{tanki}, once more
  \item \emph{Say it:} \emph{ani wa sukoshi tanki desu} — My elder brother is a little short-tempered.
  \item \emph{Recall:} say the Japanese for stubborn, then the Japanese for impatient, then say \emph{tanki} again
\end{itemize}

\begin{tcolorbox}[breakable,colback=teal!4,colframe=teal!35!black,title={Before you move on}]
What does \ja{たんき} mean? (\textquotedblleft{}Short-tempered\textquotedblright{}.) Say it once more.
\end{tcolorbox}
\section[\texorpdfstring{\ja{のんびり} (nonbiri)}{のんびり (nonbiri)}]{\ja{のんびり} (nonbiri) — leisurely; in a relaxed way}

\label{lesson:JA-C438-nonbiri}



\begin{quote}
Before the new one: say the Japanese for impatient, then the Japanese for short-tempered.
\end{quote}

\subsection*{You'll want to know: \ja{のんびり}}
\textbf{\ja{のんびり}} — \emph{nonbiri} — \textquotedblleft{}leisurely; in a relaxed way\textquotedblright{}.

\begin{quote}
\textbf{\ja{しゅうまつは} \ja{うちで} \ja{のんびりしました。}} — \emph{shūmatsu wa uchi de nonbiri shimashita} — \textquotedblleft{}I took it easy at home at the weekend.\textquotedblright{}
\end{quote}

\begin{cousinweb}
\textbf{\ja{のんびり}} (\emph{nonbiri}) is a native Japanese mimetic word, inherited rather than borrowed from Chinese, and its older origin is not settled (JMdict).
\end{cousinweb}

\begin{grammarlens}[title={What you've built}]
One more word for describing what you lived and felt.
\end{grammarlens}

\subsection*{Your turn}
\begin{itemize}
\raggedright
  \item \emph{Say it:} \emph{nonbiri}
  \item \emph{Say it:} \emph{nonbiri}, once more
  \item \emph{Say it:} \emph{shūmatsu wa uchi de nonbiri shimashita} — I took it easy at home at the weekend.
  \item \emph{Recall:} say the Japanese for impatient, then the Japanese for short-tempered, then say \emph{nonbiri} again
\end{itemize}

\begin{tcolorbox}[breakable,colback=teal!4,colframe=teal!35!black,title={Before you move on}]
What does \ja{のんびり} mean? (\textquotedblleft{}Leisurely; in a relaxed way\textquotedblright{}.) Say it once more.
\end{tcolorbox}
\section[\texorpdfstring{\ja{いばる} (ibaru)}{いばる (ibaru)}]{\ja{いばる} (ibaru) — to act big; to boss people around}

\label{lesson:JA-C438-ibaru}



\begin{quote}
Before the new one: say the Japanese for short-tempered, then the Japanese for leisurely.
\end{quote}

\subsection*{You'll want to know: \ja{いばる}}
\textbf{\ja{いばる}} — \emph{ibaru} — \textquotedblleft{}to act big; to boss people around\textquotedblright{}.

\begin{quote}
\textbf{\ja{あの} \ja{せんぱいは} \ja{いつも} \ja{いばっています。}} — \emph{ano senpai wa itsumo ibatte imasu} — \textquotedblleft{}That senior acts big all the time.\textquotedblright{}
\end{quote}

Polite \textbf{\ja{いばります}}, \emph{ibarimasu}; a godan verb. Past (polite): \textbf{\ja{しゃちょうが} \ja{いばりました。}}, \emph{shachō ga ibarimashita}.

\begin{cousinweb}
\textbf{\ja{いばる}} (\emph{ibaru}) is a native Japanese word (wago), inherited rather than borrowed from Chinese. JMdict gives two kanji spellings that differ in their first character, so the spelling was attached to the word later (JMdict).
\end{cousinweb}

\begin{grammarlens}[title={What you've built}]
One more word for describing what you lived and felt.
\end{grammarlens}

\subsection*{Your turn}
\begin{itemize}
\raggedright
  \item \emph{Say it:} \emph{ibaru}
  \item \emph{Say it:} \emph{ibaru}, once more
  \item \emph{Say it:} \emph{ano senpai wa itsumo ibatte imasu} — That senior acts big all the time.
  \item \emph{Recall:} say the Japanese for short-tempered, then the Japanese for leisurely, then say \emph{ibaru} again
\end{itemize}

\begin{tcolorbox}[breakable,colback=teal!4,colframe=teal!35!black,title={Before you move on}]
What does \ja{いばる} mean? (\textquotedblleft{}To act big; to boss people around\textquotedblright{}.) Say it once more.
\end{tcolorbox}
